\begin{corollaire}
\label{VIII.1.5}
Sea $N$ un~$A$-módulo. Entonces la sucesión de aplicaciones entre
conjuntos
$$
\xymatrix@C=.5cm{N\ar[r] & N'\ar@/^2pt/[r]\ar@/_2pt/[r] & N''}
$$
es exacta.
\end{corollaire}

Sea, en efecto, $A_{1}$ una~$A$-álgebra fielmente plana. Para demostrar
que la sucesión considerada es exacta, basta probar que lo es la sucesión
que se deduce de ella por el cambio de anillos~$A\to A_{1}$. Ahora bien,
esta última, como se ve inmediatamente, es la correspondiente al
\ifthenelse{\boolean{orig}}
{$A$-módulo}
{$A_1$-módulo}
$N_{1}=N\otimes_{A}A_{1}$ y a la $A_{1}$-álgebra
$A_{1}'=A_{1}\otimes_{A}A'$. Basta, pues, hallar un $A_{1}$ fielmente
plano sobre~$A$ tal que $\Spec(A_{1}')\to \Spec(A_{1})$ sea un morfismo
de descenso estricto para~$\cal{F}$. En efecto, basta tomar $A_{1}=A'$,
pues entonces el morfismo anterior admite un morfismo inverso por la derecha y,
por tanto, en virtud de~\cite{VIII.D}, es un morfismo de descenso efectivo
para cualquier categoría fibrada sobre~$\Sch$.

Queda, por último, demostrar que, si $N'$ es un~$A'$-módulo provisto de
un dato de descenso para~$A\to A'$, es decir, de un isomorfismo
$$
\varphi\colon N_{1}'\isomto N_{2}'
$$
entre los dos módulos deducidos de
\ifthenelse{\boolean{orig}}
{~$N$}
{~$N'$}
por los cambios de anillos
$\xymatrix@C=.5cm{A'\ar@/^2pt/[r]\ar@/_2pt/[r] & A'\otimes_{A}A'}$,
entonces
\marginpar{199}
$N'$ es isomorfo, junto con su dato de descenso, a un módulo de la forma
$N\otimes_{A}A'$. Teniendo en cuenta~\Ref{VIII.1.5}, se ve fácilmente
que este enunciado equivale al siguiente:

\begin{lemme}
\label{VIII.1.6}
Sea $N'$ un~$A'$-módulo provisto de un dato de descenso relativamente
a~$A\to A'$ (donde $A'$ es una $A$-álgebra). Sea $N$ el sub-$A$-módulo
de~$N'$ formado por los~$x$ tales que
$\varphi(x\otimes_{A'}1_{A'})=1_{A'}\otimes_{A'}x$, y consideremos el
homomorfismo canónico
$$
N\otimes_{A}A'\to N'\quoi,
$$
(que es entonces compatible con los datos de descenso). Si $A'$ es
fielmente plana sobre~$A$, este homomorfismo es un isomorfismo.
\end{lemme}

Demostremos el lema. Sea de nuevo $A_{1}$ una $A$-álgebra fielmente
plana. Para demostrar que el morfismo considerado es un isomorfismo,
basta probar que pasa a serlo después del cambio de anillos
\ifthenelse{\boolean{orig}}
{$A_{1}\to A$.}
{$A\to A_1$.}
Ahora bien, usando la planitud de~$A_{1}$ sobre~$A$, se ve que el
homomorfismo así obtenido no es otro que el que se obtendría directamente
en términos del módulo~$N'\otimes_{A}A_{1}$
sobre~$A_{1}'=A'\otimes_{A}A_{1}$, provisto del dato de descenso relativamente
a~$A_{1}\to A_{1}'$ que se deduce canónicamente, por cambio de anillos,
del que estaba dado sobre~$N'$. Así, basta hallar un~$A_{1}$ fielmente
plano sobre~$A$ tal que $\Spec(A_{1}')\to \Spec(A_{1})$ sea un morfismo
de descenso efectivo para~$\cal{F}$. Se toma entonces, como antes,
$A_{1}=A'$. Esto concluye la demostración de~\Ref{VIII.1.6} y, con ella,
la demostración de~\Ref{VIII.1.1}.

\begin{corollaire}[Descenso de secciones de módulos]
\label{VIII.1.7}
\index{descenso de secciones de módulos|hyperpage}%
Sea $g\colon S'\to S$ un morfismo de preesquemas fielmente plano y
cuasicompacto. Para todo módulo cuasicoherente~$G$ sobre~$S$, sean $G'$ y
$G''$ sus imágenes inversas sobre~$S'$ y $S''=S'\times_{S}S'$, y
consideremos el diagrama de homomorfismos de módulos sobre~$S$:
$$
\xymatrix@C=.5cm{G\ar[r] & g_*(G')\ar@/^2pt/[r]\ar@/_2pt/[r] & h_*(G'')}
$$
(donde $h\colon S''\to S$ es el morfismo estructural). Este diagrama es
\emph{exacto}.
\end{corollaire}

En
\marginpar{200}
efecto, esto significa que, para todo abierto~$U$ de~$S$, el diagrama
correspondiente formado por las secciones sobre~$U$ es exacto. Puede
suponerse, evidentemente, que $U=S$, y la exactitud en cuestión es
entonces un caso particular de~\Ref{VIII.1.2}, obtenido tomando
$F=\mathcal{O}_{S}$.

Como el funtor imagen inversa de módulos es exacto por la derecha, se
deduce formalmente de~\Ref{VIII.1.1}:

\begin{corollaire}[Descenso de módulos cociente]
\label{VIII.1.8}
\index{descenso de módulos cociente|hyperpage}%
Con las notaciones de~\Ref{VIII.1.7}, sea además, para todo módulo
cuasicoherente~$F$ sobre un preesquema, $\Quot(F)$ el conjunto de los
módulos cociente cuasicoherentes de~$F$. Con esta convención, el diagrama
de aplicaciones entre conjuntos
$$
\xymatrix@C=.5cm{\Quot(G)\ar[r] & \Quot(G')\ar@/^2pt/[r]\ar@/_2pt/[r] &
\Quot(G'')}
$$
es exacto.
\end{corollaire}

(Evidentemente, se tendría el mismo enunciado con los submódulos en lugar
de los módulos cociente, pues unos y otros se corresponden
biunívocamente). Tomando, en particular, $G=\mathcal{O}_{S}$, se obtiene:

\begin{corollaire}[Descenso de subpreesquemas cerrados]
\label{VIII.1.9}
\index{descenso de subpreesquemas cerrados|hyperpage}%
Para todo preesquema~$X$, sea $H(X)$ el conjunto de los subpreesquemas
cerrados de~$X$. Con esta notación y bajo las condiciones
de~\Ref{VIII.1.7}, el siguiente diagrama de aplicaciones entre conjuntos
$$
\xymatrix@C=.5cm{H(S)\ar[r] & H(S')\ar@/^2pt/[r]\ar@/_2pt/[r] & H(S'')}
$$
\emph{es exacto}.
\end{corollaire}

Conviene completar el teorema~\Ref{VIII.1.1} con el resultado siguiente:

\begin{proposition}[Descenso de propiedades de módulos]
\label{VIII.1.10}
\index{descenso de propiedades de módulos|hyperpage}%
Sean $g\colon S'\to S$ un morfismo fielmente plano y cuasicompacto, y
$F$ un módulo cuasicoherente sobre~$S$. Para que~$F$ sea de tipo finito,
\resp de presentación finita, \resp localmente libre y de tipo finito,
es necesario y suficiente que su imagen inversa $F'$ sobre~$S'$ lo sea.
\end{proposition}

Solo hay que demostrar la suficiencia. Puede suponerse, evidentemente,
que~$S$ es afín,
\marginpar{201}
y, reemplazando entonces~$S'$ por
\ifthenelse{\boolean{orig}}
{la}
{una}
suma de abiertos afines que recubren~$S'$, se queda reducido al caso en
que~$S'$ es también afín. Nuestro enunciado equivale entonces al
siguiente:

\begin{corollaire}
\label{VIII.1.11}
Sean $A$ un anillo, $A'$ una~$A$-álgebra fielmente plana, $M$
un~$A$-módulo y $M'$ el~$A'$-módulo $M\otimes_{A}A'$. Para que~$M$ sea de
tipo finito (\resp de presentación finita, \resp localmente libre de tipo
finito), es necesario y suficiente que~$M'$ lo sea.
\end{corollaire}

En efecto, se tiene $M=\varinjlim_i M_{i}$, donde los~$M_{i}$ son los
submódulos de tipo finito de~$M$. Por consiguiente,
$M'=\varinjlim_i M_{i}'$, y, si $M'$ es de tipo finito, $M'$ es igual a
uno de los~$M_{i}'$; por planitud fiel, $M$ es entonces igual a~$M_{i}$ y,
por tanto, $M$ es de tipo finito. En consecuencia, existe una sucesión
exacta
$$
0\to R\to L\to M\to 0\quoi,
$$
con $L$ libre de tipo finito, de donde se obtiene una sucesión exacta
$$
0\to R'\to L'\to M'\to 0\quoi,
$$
con $L'$ libre de tipo finito. Así pues, si $M'$ es de presentación
finita, $R'$ es de tipo finito; por lo que precede, $R$ es de tipo finito
y, por tanto, $M$ es de presentación finita. Por último, decir que~$M$ es
localmente libre y de tipo finito significa que es de presentación
finita y plano (\cf \Ref{IV} en el caso noetheriano; el caso general se
deja al lector). Como cada una de estas propiedades desciende, lo mismo
ocurre con su conjunción, lo que concluye la demostración.

\begin{remarque}
\label{VIII.1.12}
La conjunción de~\Ref{VIII.1.1} y~\Ref{VIII.1.10} muestra que el
enunciado~\Ref{VIII.1.1} sigue siendo válido cuando se reemplaza la
categoría fibrada~$\cal{F}$ por la subcategoría fibrada formada por los
módulos cuasicoherentes de tipo finito, \resp de presentación finita,
\resp localmente libres de tipo finito, \resp localmente libres de rango
dado~$n$.
\end{remarque}
